\documentclass[12pt]{article}
\usepackage[utf8]{inputenc}
\usepackage[T1]{fontenc}
\usepackage{hyperref}
\title{Swedish Sign Language Language Resources for K-12 Curriculum Development}
\author{Compiled by Claude Sonnet 5.5 (claude-sonnet-5-5) for Andr\'as Kornai}
\date{}
\begin{document}
\maketitle

\section{Introduction}
Swedish Sign Language (svenskt teckenspråk, ISO 639-3 swl, Glottocode swed1236) is the language of Sweden's Deaf community. Sweden has an established system for teaching it, through the special schools run by the Swedish National Agency for Special Needs Education and Schools (SPSM) and through a hearing-learner subject in the compulsory-school regulations. The resources below come from SPSM, Skolverket and the Swedish Deaf Association (SDR).

\section{Primary Educational Resources}
\begin{itemize}
\item SPSM framework syllabus (ramkursplan) in Swedish Sign Language, grades 1 to upper secondary. It has no grading criteria and targets, among others, hearing children of deaf parents (CODA), siblings of deaf children and part-time pupils: \url{https://www.spsm.se/contentassets/7e31ce7c5ea7406cbb94c7e5b29b3ff0/ramkursplan-plan.pdf}
\item Skolverket regulation SKOLFS 2023:275, which includes a subject of sign language for hearing pupils (teckenspråk för hörande): \url{https://skolfs.skolverket.se/api/document/ANDRINGSFORFATTNING/2023:275/pdf}
\item SPSM information sheet on Vänerskolan, a special school; it describes the aim of developing two full languages, Swedish Sign Language and Swedish: \url{https://spsm.se/siteassets/vara-skolor-2025/informationsblad-alla-skolor-2425/infoblad_vanerskolan.pdf}
\end{itemize}

\section{Linguistic and Reference Materials}
\begin{itemize}
\item SPSM publication on learning, developing and using sign language: \url{https://webbutiken.spsm.se/globalassets/publikationer/filer/lara-utveckla-och-anvanda-teckensprak-tillganglig-version.pdf/}
\item SPSM publication on sign language, spoken and written language: \url{https://webbutiken.spsm.se/globalassets/publikationer/filer/teckensprak-talat-och-skrivet-sprak-pdf}
\end{itemize}

\section{Religious and Community Resources}
\begin{itemize}
\item Swedish Deaf Association (SDR) education policy document: \url{https://sdr.org/wp-content/uploads/2023/08/utbildningspolicy.pdf}
\end{itemize}

\section{Digital and Audio Resources}
The searches did not return sign-language video lexicons or video corpora. Video resources (for example, signed lexicons) should be looked up separately before use.

\section{Curriculum Development Context}
Skolverket has started work on new syllabi in Swedish Sign Language. SPSM carried out a preparatory assignment in autumn 2025 on how special-school teachers experience the current syllabi. The report is dated November 2025 and feeds into Skolverket's further work: \url{https://www.skolverket.se/download/18.2cb418a019afcca00f8eb2f/1765373625660/F%C3%B6rberedande%20uppdrag%20f%C3%B6r%20kursplaner%20i%20teckenspr%C3%A5k.pdf}

\end{document}